\documentclass[11pt]{article}

\usepackage[margin=1in]{geometry}
\usepackage{amsmath, amssymb}
\usepackage{booktabs}

\setlength{\parindent}{0pt}
\setlength{\parskip}{4pt}

\newcommand{\code}[1]{\texttt{#1}}

\title{The GP Fit: Model, Priors, Fitting, and What Comes Out}
\author{\texttt{GenerateTable\_CGW}}
\date{2026 November 3}

\begin{document}
\maketitle

\begin{abstract}
Reference for the Gaussian-process stage of the pipeline (\code{gp\_pipeline/}): the kernel,
how components are seeded from the periodogram, the priors, the staged optimisation, the
columns it writes, and the quality flags attached to them. Validation numbers quoted here
come from the synthetic injection--recovery set; the derivations and the history of the
red-noise term live in \code{docs/fixing\_RN\_fitting.tex}, and the peak-finding statistics
in \code{docs/peak\_finders.tex}.
\end{abstract}

% ===========================================================================
\section{What the GP is for}
% ===========================================================================

Each row of the table is one light curve: a cluster observed in one sector combination,
resampled to 30-min bins and normalised per sector. The GP models it as a sum of
quasi-periodic oscillators sitting on a red-noise continuum, plus white noise:
\begin{equation}
k(\tau) = \underbrace{\sum_{i=1}^{m} k_{\rm SHO}(\tau;\sigma_i, P_i, Q_i)}_{\text{signals}}
\;+\; \underbrace{k_{\rm RN}(\tau; S_0, \alpha)}_{\text{continuum}}
\;+\; \underbrace{(\sigma_{\rm phot}^2 + \sigma_{\rm jit}^2)\,\delta_\tau}_{\text{white noise}} .
\end{equation}
All terms are \code{tinygp} quasiseparable kernels, so the likelihood costs $O(N)$ rather
than $O(N^3)$ and a 1300-point row fits in tens of seconds.

The red-noise term earns its place: without it the SHO components drift into modelling the
continuum, which corrupts the periods that the science actually wants. That failure is what
the $Q$ prior (\S4) and the red-noise term exist to prevent.

% ===========================================================================
\section{The kernel}
% ===========================================================================

\subsection{Signal components}

Each is a simple harmonic oscillator, \code{quasisep.SHO}, free in amplitude $\sigma_i$,
period $P_i$ and quality factor $Q_i$. Its PSD peaks at $1/P_i$ with width set by $Q_i$:
a coherent sinusoid drives $Q$ to the cap, spot modulation with a finite lifetime does not.
Periods are bounded to the search band (0.083--10\,d).

\subsubsection{Bounds on $Q$: a decay-time cap, not a $Q$ cap}
\label{sec:qbounds}

For $Q > 1/2$ the SHO kernel is an envelope times a cosine,
$k(\tau) \propto e^{-\omega_0 \tau / 2Q}\cos(\ldots)$ (Foreman-Mackey et al.\ 2017, eq.~23),
so the amplitude e-folds on the \textbf{decay time}
\begin{equation}
\tau_d = \frac{2Q}{\omega_0} = \frac{Q\,P}{\pi},
\qquad\text{i.e.}\qquad Q = \pi\,\frac{\tau_d}{P},
\end{equation}
which is the $(\rho, \tau)$ parameterisation \code{celerite2} offers for \code{SHOTerm}.
$Q/\pi$ is the number of coherent cycles before the signal forgets its phase.

\textbf{Upper bound.} Once $\tau_d$ exceeds the length of the light curve the likelihood is
flat in $Q$: a component coherent for 27\,d and one coherent for ever are the same thing in
27\,d of data, and the optimiser drifts to whatever edge it was given. Nicholson \& Aigrain
(2022) show, on spot-model simulations, that the recovered evolution timescale of a
quasi-periodic GP tracks the truth only up to about half the baseline and scatters freely
beyond it. The cap is therefore on $\tau_d$, and in $Q$ it is \emph{period-dependent}:
\begin{equation}
\tau_d \le \tau_{\rm cap}
\quad\Longleftrightarrow\quad
Q \le Q_{\max}(P) = \frac{\pi\,\tau_{\rm cap}}{P},
\qquad \tau_{\rm cap} = \min(27\,\mathrm{d},\ \text{baseline}).
\end{equation}
(The frequency-domain statement is the same: the SHO line, FWHM $\simeq f_0/Q$, must not be
narrower than the Rayleigh resolution $1/T$, which gives $Q < T/P$ --- the same scaling, a
factor $\pi$ tighter.) At the defaults this is $Q_{\max} = 1020$ at $P = 0.083$\,d, 85 at
1\,d, 17 at 5\,d and 8.5 at 10\,d.

$\tau_{\rm cap}$ is pinned to \textbf{one TESS sector for every row}, including multi-sector
combinations. A stitching offset between sectors can only be absorbed as an 80-day coherent
signal if the box allows an 80-day decay time, and the $Q$ feature should mean the same thing
whatever the sector count. Rows shorter than a sector use their own baseline.

Before 2026-11 the box was a flat $Q \in [1/\sqrt2, 500]$. That equals the baseline cap at
exactly one period (0.17\,d for a sector), so below it the box cut off decay times the data
could resolve, and above it the box allowed decay times far beyond the data: on the
2026-09-21 real table 38\% of components sat on the 500 cap, 43\% had $\tau_d$ longer than
their own light curve, and above $P = 2$\,d the \emph{median} component already exceeded
$\pi T/P$. The new cap \emph{loosens} the bound below $\sim 0.5$\,d (where the pile-up at 500
was mostly the flat cap biting) and \emph{tightens} it above $\sim 2$\,d. A component at the
cap now reads ``coherent for the whole window'', which is at least an interpretable
statement; it is still not a measurement.

\textbf{Lower bound.} Physical rather than arbitrary: an SHO's PSD has a peak at non-zero
frequency only for $Q > 1/\sqrt2$, so below that the component is a continuum, not an
oscillator --- the same value the old fixed-$Q$ red-noise SHO used. It is a hard bound. The
soft prior (\S4) does the real work above it.

\textbf{How the bounds are enforced.} L-BFGS-B takes only box bounds fixed before the fit,
and both limits on $Q$ move with the period. So the parameter vector carries, per
component, a \emph{coherence fraction} $s \in [0,1]$ instead of $\log Q$, and
\begin{equation}
\log Q = \log Q_{\min} + s\,\bigl[\log Q_{\max}(P) - \log Q_{\min}\bigr]
\end{equation}
is evaluated inside \code{unpack\_theta} from the current period
(\code{gp\_fit.q\_from\_coherence}). The box on $s$ is $[0,1]$, so both limits are exact hard
walls; $s = 0$ is the continuum limit and $s = 1$ is ``coherent for the whole cap''. The
kernel, the soft prior and every reported feature use the derived $Q$; nothing downstream
sees $s$ except the \code{gp\_sho\_coherence} column, which is the cleanest ``is it at the
cap?'' diagnostic. Fits pickled before the change carry no \code{tau\_cap} key and unpack in
the old $\log Q$ layout (\code{gp\_fit.q\_layout\_kwargs}), so old tables still load.

\subsection{Red noise: a power law built from OU terms}

A single fixed-$Q$ SHO cannot represent $S(f) \propto f^{-\alpha}$ --- it is flat below its
break and $f^{-4}$ above, so its effective slope carries no information about $\alpha$
(measured correlation with truth: $-0.07$). Instead the continuum is a fixed basis of $N=9$
Ornstein--Uhlenbeck (\code{quasisep.Exp}) terms at log-spaced frequencies $f_j$ spanning the
band plus 1.2 decades of padding at each end, with weights tied to a single slope:
\begin{equation}
S_{\rm RN}(f) = \sum_j \frac{c_j}{1 + (f/f_j)^2},
\qquad c_j = S_0\,f_0^{\alpha} f_j^{-\alpha} e^{u_j(\alpha)},
\qquad \sigma_j^2 = \tfrac{\pi}{2} c_j f_j .
\end{equation}
Only $S_0$ (the PSD at the pivot $f_0 = 1$\,d$^{-1}$) and $\alpha$ are free, whatever $N$ is,
so the parameter count matches the old SHO term and BIC comparisons still line up.

The corrections $u_j(\alpha)$ carry the \textbf{exposure filter}. Every flux value is an
average over its 30-min bin, so the data hold $S(f)\,\mathrm{sinc}^2(fE)$, not $S(f)$. A
kernel that ignores this fits $\alpha$ about 0.3 too steep. \code{exposure\_weight\_table}
solves once per (basis, exposure) for the $u_j$ that reproduce the filtered power law, and
\code{jnp.interp} reads them at the current $\alpha$ inside the likelihood. Accuracy is
0.02--0.05\,dex over 0.1--12\,d$^{-1}$; the kernel stays a sum of OU terms, so fits cost the
same. $S_0$ and $\alpha$ then describe the \emph{intrinsic} spectrum.

\textbf{Limits.} A sum of OU terms cannot fall faster than $f^{-2}$ anywhere, so
$\alpha$ is bounded to $[0.05, 1.95]$ and the in-band slope saturates near 1.8. On real data
only 1--2\% of rows have a periodogram slope steeper than 1.8, so this rarely binds; a
Mat\'ern-3/2 basis would lift it if that changes.

% ===========================================================================
\section{Seeding: which peaks become components}
% ===========================================================================

The GP does not search for periods. It is seeded from the stored periodogram, and a bad seed
is never recovered, so seeding matters as much as the fit.

\code{find\_all\_significant\_peaks} walks the periodogram greedily, taking the largest
excess over the red-noise continuum $C(f)$, masking $\pm5\%$ in frequency around it, and
repeating while any bin stays above the significance threshold $\gamma C(f)$ (Vaughan 1\%).
\code{threshold\_mode='red\_noise'} is the pipeline default; \code{'red\_noise\_subtracted'}
orders candidates by $P - C$ rather than $P - \gamma C$, which is unbiased in period
(\code{peak\_finders.tex}).

$C(f)$ comes from the \textbf{peak-masked} continuum fit
(\code{lc\_lsp.\_fit\_red\_noise\_vaughan\_masked}). Fitting every bin lets the signals bias
the continuum steep --- median $\alpha$ error $+0.71$ with 7--10 injected oscillators --- which
lifts the threshold at low frequency and hides real peaks. Masking cut the
detectable-but-missed oscillator count from 157 to 42 on the synthetic set.

\textbf{Harmonics and sidelobes are not suppressed} (\code{harmonic\_masking=False}): a
harmonic is a real feature of the light curve. The cost is that one physical signal can
produce several seeds --- in crowded rows only about a third of the peaks found are distinct
injections, a third are harmonics and a third are sidelobes of a peak already found. Hence
\S6.

% ===========================================================================
\section{Priors}
% ===========================================================================

All three are smooth penalties added to the negative log-likelihood, resolved at Python
level into one objective (so JIT sees a fixed structure). BIC uses the bare log-likelihood,
not the penalised objective.

\begin{description}
\item[Period (log-normal, \code{use\_lognormal\_period\_prior})] keeps each component near its
  seed, $0.5\,[(\ln P_i - \ln P_i^{\rm seed})/\sigma_{\rm ln}]^2$ with $\sigma = 0.3$ dex by
  default. Without it components wander off their peak during the free pass.
\item[$Q$ floor (\code{use\_q\_prior})] penalises $Q < \pi$ quadratically in $\log Q$
  (\code{q\_prior\_threshold}). $Q = \pi$ is a decay time of exactly one period
  (\S\ref{sec:qbounds}): Angus et al.\ (2018) argue that the evolution timescale of a
  rotating spotted star should exceed its period, $l > P$, and the spot simulations behind
  Aigrain et al.\ (2015) and Nicholson \& Aigrain (2022) use lifetimes of $1$--$10\,P$, i.e.\
  $Q \approx 3$--$30$. A low-$Q$ SHO is broad, and a broad component is a continuum in
  disguise; this is what stops signal components from eating the red noise. The threshold
  was 2 ($\tau_d = 0.64\,P$) before 2026-11. The prior acts on the \emph{derived} $Q$, so it
  is unchanged by the coherence-fraction parameterisation.
\item[Red-noise slope in the white-noise limit] --- see \S5.
\end{description}

% ===========================================================================
\section{The white-noise limit}
% ===========================================================================

White noise is the $\alpha \to 0$ limit of red noise, but the kernel also has a jitter term,
and a flat red-noise component is degenerate with it. Given a light curve with no red noise
the fit therefore zeroes the red-noise \emph{amplitude} and leaves $\alpha$ floating --- on the
first real run, 33\% of rows ended that way with a median $\alpha$ of 1.76 and half piled at
the upper bound, an artefact that correlates with amplitude and could easily be mistaken for
signal.

\code{\_rn\_alpha\_wn\_neg\_log\_prior} adds $w_{\rm pr}\alpha^2\,\sigma(-\log_{10}(r/r_0)/\Delta)$
with $r = S_0 / \mathrm{WN\ PSD}$, defaults $w_{\rm pr}=1$, $r_0 = 0.1$, $\Delta = 0.5$
decades. Where red noise is real the weight is negligible (0.04 at $r=4$); where it is
absent, $\alpha$ is pulled to the flat limit. On the synthetic set the flagged rows moved
from a median $\alpha$ of 1.04 (35\% at the bound) to 0.05 (4\% at the bound), while rows with
real red noise were unchanged: median error $0.087 \to 0.078$, $r$ with truth
$0.798 \to 0.828$.

\textbf{This is an imposed value, not a measurement.} Report those rows as "consistent with
flat" and keep \code{gp\_rn\_wn\_ratio} beside them.

% ===========================================================================
\section{Component budget and pruning}
% ===========================================================================

Because harmonics and sidelobes each produce seeds, \code{max\_components} is spent on
duplicates in crowded rows. The GP handles a duplicate gracefully --- it is fitted at a median
\textbf{4\%} of the amplitude of the component holding the real signal, at $Q \sim 40$ against
$Q \sim 500$ for a real match --- so duplicates waste budget rather than corrupt the fit.

Two changes follow, and they work together:

\begin{center}
\begin{tabular}{lcc}
\toprule
 & \code{max\_components}=10 & =16 \\
\midrule
rows at the cap (of 246)                       & 69 & \textbf{39} \\
rows with neither quality flag                 & 160 & \textbf{187} \\
$\alpha$ error, all rows with red noise        & 0.097 & \textbf{0.091} \\
raw purity / completeness                      & 0.736 / 0.928 & 0.665 / \textbf{0.974} \\
filtered ($\sigma \geq 10\%$, $Q \geq 5$)      & 0.938 / 0.846 & 0.964 / \textbf{0.887} \\
mean fit time per row                          & 33\,s & 41\,s \\
\bottomrule
\end{tabular}
\end{center}

At matched purity the larger budget buys 3.8 points of completeness for 24\% more fit time.
Note that BIC does \emph{not} prefer the bigger models (median $\Delta$BIC $= 0$, better in
4\% of rows): the extra components are redundant, and the point is only that duplicates stop
crowding out real peaks.

\code{gp\_table\_io.filter\_gp\_components(table)} does the pruning afterwards, writing
\code{gp\_periods\_filtered} and friends; \code{synth/compare.py:component\_filter\_scan}
measures what a cut costs. Keeping this post hoc means the trade stays a knob rather than a
decision baked into the fit. \textbf{Caveat:} the injections are pure sinusoids, driven to
the $Q$ cap, so the $Q$ separation is sharper on synthetic data than it will be for
quasi-periodic spot modulation. The default was therefore $Q_{\min} = 5$, not the 20 that the
synthetic numbers alone would suggest: on the real table $Q_{\min} = 20$ discarded obvious
signals --- NGC2136 shows a clear 4.4\,d modulation by eye, fitted with three components at
$Q \approx 14$, all of which that cut removed. Real components sit in $Q = 2$--20 far more
often than synthetic ones (24\% vs 17\%), which is exactly what finite coherence looks like.

\textbf{The cut is now in cycles.} With the period-dependent cap of \S\ref{sec:qbounds} a
flat $Q \ge 5$ means different things at different periods: at $P = 10$\,d the ceiling is
$Q = 8.5$, so the cut would sit almost on it, while at $P = 0.1$\,d it is 1\% of the ceiling.
\code{filter\_gp\_components} therefore cuts on $\tau_d / P = Q/\pi \ge$
\code{n\_cycles\_min}, default 1.6, which is exactly the old $Q \ge 5$; \code{Q\_min} is
still accepted and overrides it. The table above compares component budgets under the old
flat cap; the filter scan under the decay-time cap (same 246 rows, 1632 components, 611
detectable injections) is:

\begin{center}
\begin{tabular}{lccc}
\toprule
cut & components & purity & completeness \\
\midrule
none                                    & 1632 & 0.684 & 0.977 \\
$\sigma \ge 10\%$                       &  559 & 0.961 & 0.889 \\
cycles $\ge 1.6$ ($Q \ge 5$)            & 1529 & 0.708 & 0.972 \\
$\sigma \ge 10\%$ and cycles $\ge 1.6$  &  537 & \textbf{0.983} & \textbf{0.887} \\
$\sigma \ge 10\%$ and cycles $\ge 6.4$ ($Q \ge 20$)  & 316 & 0.987 & 0.527 \\
$\sigma \ge 10\%$ and cycles $\ge 32$ ($Q \ge 100$)  &  50 & 1.000 & 0.082 \\
\bottomrule
\end{tabular}
\end{center}

Two things changed. The amplitude cut does most of the work, as before, and the default
coherence cut now adds purity for free (0.961 $\to$ 0.983 at the same completeness; under
the flat cap it was 0.964). But \textbf{a $Q \ge 20$ cut is no longer a tightening, it is a
period cut}: the ceiling $\pi\,27/P$ falls below 20 for $P > 4.2$\,d, so every long-period
component is removed whatever its coherence, and completeness halves. Any cut above
$\sim 2$ cycles must be made in \code{gp\_sho\_coherence}, not in $Q$.

The remark above that real components sit in $Q = 2$--20 ``far more often than synthetic
ones'' was measured under the flat cap, where $Q = 2$--20 meant finite coherence. Under the
decay-time cap that range also contains the \emph{cap itself} for $P > 4$\,d: on the
synthetic set the fraction of components in $Q = 2$--20 went from 16\% to 49\% with no
change in the injections. Read coherence from \code{gp\_sho\_coherence} (or cycles against
the cap), not from raw $Q$.

% ===========================================================================
\section{Fitting}
% ===========================================================================

L-BFGS-B with JAX gradients, in stages, each warm-started from the previous one and kept only
if it does not worsen the objective:
\begin{enumerate}
\item \textbf{periods frozen} at their seeds, so $\sigma$, $Q$, the red-noise term and jitter
  find their level while the period is pinned;
\item \textbf{everything free}.
\end{enumerate}
With \code{fit\_full\_only=False} the search runs $m = 1 \ldots n$ components and keeps the
best BIC; the pipeline sets \code{fit\_full\_only=True} and fits the full seed list once.

\textbf{The optimisation is multimodal} and this is not fully solved. On real rows the free
fit sometimes ends below a fit with $\alpha$ held at 1 (4/30 random rows, 13/30 steep rows),
which cannot happen at a true optimum. Seeding $\alpha$ from the periodogram slope did not
help; releasing $\alpha$ in a final stage
(\code{rn\_alpha\_release\_last}) guarantees no-worse-than-fixed but leaves $\alpha$ stuck at
its seed in 40\% of rows and was worse overall, so it is off by default. Running both and
keeping the better objective would cost roughly twice the fit time and has not been tested.

\subsection{When no peak is significant}

If the peak search returns nothing, there is nothing to seed. The fit then uses
\textbf{zero SHO components} --- red noise plus jitter only --- rather than inventing one.

This replaced a fallback that fitted a single SHO seeded at $P=1$\,d, and the change is not
cosmetic. That filler component landed at the $Q$ bound and \emph{absorbed the continuum}:
across the 300 affected rows of the real table (24\% of it), the median \code{gp\_rn\_alpha}
read 0.05 against 0.96 for rows that found peaks, and the median red-noise-to-white-noise
power ratio read $4\times10^{-6}$. After the change those rows report $\alpha = 1.07$ and a
ratio of 0.98, i.e. ordinary red noise, and 20 units of BIC are recovered per row while the
log-likelihood falls by only 0.5 --- the discarded component was contributing nothing but
parameters. The share of rows flagged \code{gp\_rn\_near\_white\_noise} fell from 32\% to
23\%, because two thirds of that flag was this artefact rather than genuine white noise.

Rows fitted this way have \code{gp\_n\_components = 0} and empty component lists; nothing
downstream needs special handling. Without a red-noise term in the kernel the old fallback is
kept, since the kernel would otherwise be empty.

% ===========================================================================
\section{What it writes}
% ===========================================================================

\begin{center}
\begin{tabular}{ll}
\toprule
column & meaning \\
\midrule
\code{gp\_n\_components}, \code{gp\_bic}, \code{gp\_log\_likelihood} & model size and fit quality \\
\code{gp\_periods}, \code{gp\_sho\_sigmas}, \code{gp\_sho\_omegas}, \code{gp\_sho\_Qs} & per-component parameters \\
\code{gp\_sho\_taus} & per-component decay time $\tau_d = QP/\pi$ in days; at
   \code{gp\_sho\_tau\_cap} it means ``coherent for the whole window'', not a measurement \\
\code{gp\_sho\_coherence} & per-component $s \in [0,1]$, where $Q$ sits between its floor and
   its period-dependent ceiling in log; $s=1$ is at the cap (\S\ref{sec:qbounds}) \\
\code{gp\_sho\_tau\_cap} & the decay-time cap applied to the row, $\min(27\,\mathrm{d}, T)$;
   NaN for fits made under the old flat $Q$ box \\
\code{gp\_jitter} & fitted white-noise term \\
\code{gp\_rn\_alpha} & fitted intrinsic red-noise slope \\
\code{gp\_rn\_alpha\_band} & in-band slope of the fitted spectrum --- \emph{compare this to measured slopes} \\
\code{gp\_rn\_log10\_S0}, \code{gp\_rn\_rms\_band} & red-noise level at 1\,d$^{-1}$; in-band RMS \\
\code{gp\_rn\_wn\_ratio} & red-noise / white-noise power at 1\,d$^{-1}$ \\
\code{gp\_rn\_exposure\_days}, \code{gp\_rn\_basis\_freqs} & exposure filter and basis actually used \\
\code{GP\_freq}, \code{GP\_PSD}, \code{gp\_kspace\_*} & analytic PSD of the fitted kernel, and band statistics \\
\code{gp\_initial\_lsp\_peak\_periods}, \code{gp\_n\_peaks\_found} & seeds used, peaks available \\
\code{gp\_rn\_near\_white\_noise} & \textbf{flag}: no detectable red noise, $\alpha$ undefined \\
\code{gp\_components\_capped} & \textbf{flag}: more peaks than \code{max\_components} \\
\code{gp\_*\_filtered}, \code{gp\_component\_kept} & written by \code{filter\_gp\_components};
   the second is a per-component boolean aligned with \code{gp\_periods}, so a plot can draw
   kept and cut components differently \\
\bottomrule
\end{tabular}
\end{center}

\code{gp\_rn\_sigma} / \code{gp\_rn\_period} / \code{gp\_rn\_Q} appear only under the legacy
\code{rn\_model='sho'} and are NaN otherwise.

\subsection{How to use the flags}

\code{gp\_rn\_near\_white\_noise} is a hard exclusion for $\alpha$: that row has no slope to
measure, and the reported value is the prior's. \code{gp\_components\_capped} is softer ---
those rows fit $\alpha$ nearly as well as clean rows (median error 0.12 vs 0.09) and mostly
signal a crowded periodogram. \textbf{Do not drop flagged rows from a population analysis}:
the capped fraction correlates with age ($\rho = -0.18$ on the real table), so cutting on it
selects on the target.

% ===========================================================================
\section{Validation summary}
% ===========================================================================

From the synthetic injection--recovery set (246 rows, exposure-averaged sampling, decay-time
cap of 27\,d, soft floor $Q = \pi$; figures in \code{sytheticdata\_generate\_compare/figs/},
regenerated by \code{make\_figs.py} from
\code{data/cluster\_table\_Nov03\_3sectorcombos\_synthetic\_expavg\_gp\_taucap.pkl}):

\begin{center}
\begin{tabular}{lcc}
\toprule
quantity & result & note \\
\midrule
red-noise slope, GP            & bias $+0.07$, median error 0.09, $r = 0.80$ & clean rows (187 of 246) \\
red-noise slope, periodogram   & bias $+0.01$, median error 0.17, $r = 0.60$ & same rows \\
red-noise amplitude, GP        & Spearman 0.82 & in-band vs full-baseline: offset by construction \\
detectable oscillators found   & \textbf{97.7\%} & 597 of 611, 14 missed \\
component purity               & 68.4\% raw, \textbf{98.3\%} filtered & completeness 88.7\% after the cut \\
rows fitted with no components & 57 of 246 & no significant peak; red noise + jitter only \\
\bottomrule
\end{tabular}
\end{center}

\subsection{Before and after the $Q$ bound change}

The same base table was fitted with the flat box ($Q \in [1/\sqrt2, 500]$, soft floor
$Q = 2$) and with the decay-time cap (\S\ref{sec:qbounds}); the seeds are identical, so
both runs fit the same 1632 components. Figure
\code{figs/q\_bounds\_before\_after.png}; numbers in \code{figs/headline\_before\_after.csv}.

\begin{center}
\begin{tabular}{lcc}
\toprule
 & flat $Q$ box & decay-time cap \\
\midrule
detectable oscillators recovered (of 611)   & 601 (98.4\%) & 597 (97.7\%) \\
purity, raw / filtered                      & 0.684 / 0.964 & 0.684 / \textbf{0.983} \\
completeness after filter                   & 0.887 & 0.887 \\
delivered components at an unrelated frequency & 12 & \textbf{7} \\
$\alpha$ (GP), clean rows: bias / median error / $r$ & $+0.09$ / 0.09 / 0.82 & $+0.07$ / 0.09 / 0.80 \\
rows flagged no red noise / capped          & 34 / 25 & 32 / 27 \\
components on the $Q$ cap                   & 459 at $Q = 500$ (28\%) & 1332 at $s > 0.99$ (82\%) \\
median $Q$, matched / spurious              & 159 / 26 & 21 / 19 \\
median cycles $Q/\pi$, matched / spurious   & 50 / 8.4 & 6.7 / 5.9 \\
\bottomrule
\end{tabular}
\end{center}

Recovery, completeness and the red-noise parameters are unchanged within noise; the four
net losses are period swaps between neighbouring seeds in crowded 10--16-component rows,
and recovery against injected period (panel d of the figure) is flat before and after. The
delivered list is purer because the components the cap now pins are pinned at a
\emph{period-aware} ceiling, which separates real from duplicate components better than a
single number did.

The pile-up is the expected one. The injections are pure sinusoids, coherent for ever, so
under a cap that means ``coherent for the whole sector'' they must sit on it: 90\% of
matched components have $s > 0.99$, against 64\% of spurious ones. The raw-$Q$ separation
between the two populations has therefore collapsed (medians 21 against 19, where the flat
cap gave 159 against 26), and \textbf{on synthetic data $Q$ no longer discriminates; $s$ and
the amplitude do}. On real spot-modulated light curves, where decay times are finite, $s$
is expected to spread below 1 and the cycles cut to regain its meaning; that has not been
measured yet (the real table has not been refitted).

Against age on the real table, the red-noise \emph{amplitude} correlates
($\rho = -0.31$ at cluster level) while the \emph{slope} does not ($\rho \approx 0$), so
$\alpha$ is currently a diagnostic rather than a feature.

\textbf{Known open items.} The optimiser multimodality above; the $\alpha < 2$ ceiling; the
crowded shallow-$\alpha$ rows where leftover oscillator power still biases the continuum; and
the fact that the synthetic injections are pure sinusoids, which flatters both the $Q$
separation and the period recovery relative to real quasi-periodic variability.

\textbf{After the $Q$ bound change (2026-11).} The synthetic set has been refitted and
every number above is from the decay-time-capped fit except where a table says otherwise
(the \code{max\_components} comparison in \S6 and the real-table statistics in
\S\ref{sec:qbounds}). The real table has \emph{not} been refitted: the
\code{gp\_sho\_Qs} distribution, the $Q = 2$--20 fractions and the NGC2136 example quoted
in \S6 all come from the flat-cap fit, and the component filter on real data should be
re-checked once it is.

\end{document}
